\documentclass[11pt,a4paper]{article}
\usepackage[utf8]{inputenc}
\usepackage[T1]{fontenc}
\usepackage[brazil]{babel}
\usepackage{lmodern}
\usepackage[margin=2.5cm]{geometry}
\usepackage{booktabs}
\usepackage{adjustbox}
\usepackage{array}
\usepackage{xcolor}
\usepackage{enumitem}
\usepackage{hyperref}
\setlist{itemsep=2pt}

\newcommand{\todo}[1]{\textcolor{red}{[\textbf{A FAZER:} #1]}}
% caixa de "ficha" de cada tabela de dados
\newcommand{\ficha}[4]{%
\begin{center}\small
\begin{tabular}{>{\bfseries}p{3.2cm}p{11.5cm}}
\toprule
Tabela & #1 \\
Cada linha é & #2 \\
Quantas linhas & #3 \\
De onde vem & #4 \\
\bottomrule
\end{tabular}
\end{center}}

\title{Justiça e IA: um juiz que segue regras pode decidir de forma enviesada?}
\author{}
\date{}

\begin{document}
\maketitle

% =====================================================================
\subsection{Metodologia}
% =====================================================================

\subsubsection*{A pergunta}

Se escrevermos a lei como um programa de computador (um ``juiz virtual'' que só segue
regras), as decisões ficam livres de viés? A hipótese do trabalho é que não: para
transformar a lei em regras, alguém precisa preencher o que a lei deixa em aberto, e
essas escolhas mudam o resultado.

\subsubsection*{O caso estudado: o tráfico privilegiado}

O crime de tráfico (art.~33 da Lei de Drogas, Lei 11.343/2006) tem pena de 5 a 15 anos.
O parágrafo 4º desse artigo permite \textbf{reduzir a pena de 1/6 a 2/3} quando o réu
cumpre quatro condições ao mesmo tempo:
\begin{enumerate}
  \item é \textbf{primário} (não é reincidente);
  \item tem \textbf{bons antecedentes};
  \item \textbf{não se dedica a atividades criminosas};
  \item \textbf{não faz parte de organização criminosa}.
\end{enumerate}
Com a redução máxima, uma pena de 5 anos cai para 1 ano e 8 meses, em regime aberto.
Sem a redução, fica em 5 anos ou mais, em regime semiaberto ou fechado.
Escolhemos esse ponto porque é o pedido mais comum nos recursos e porque a lei
\textbf{não diz} duas coisas: o que significa ``dedicar-se a atividades criminosas'' e
como escolher a fração entre 1/6 e 2/3 (manda só ``considerar a natureza e a quantidade
da droga'').

\subsubsection*{De onde vieram os dados}

\begin{itemize}
  \item \textbf{Decisões do Superior Tribunal de Justiça (STJ)}, publicadas pelo próprio
        tribunal em arquivos abertos, um por mês. Usamos as duas turmas que julgam
        crimes (5ª e 6ª Turma), de junho a agosto de 2026.
  \item \textbf{A ficha de cada processo} no site do STJ: quem são as partes, os recursos
        apresentados e o texto completo das decisões (em PDF).
  \item \textbf{O texto da lei} em vigor (Lei de Drogas e Código Penal).
\end{itemize}

Ficaram de fora os casos que não tratam de uma condenação por tráfico: prisões antes da
sentença, execução da pena, atos de adolescentes e o crime de associação sem tráfico.

\subsubsection*{As tabelas de dados}

O trabalho usa quatro tabelas. Para cada uma, dizemos o que é uma linha, quantas
linhas há e quais colunas (variáveis) ela tem.

\paragraph{Tabela 1 -- Decisões do STJ (base grande).}
\ficha{\texttt{dados/ementas\_t5.csv} e \texttt{dados/ementas\_t6.csv}}
      {um processo de tráfico, na decisão mais recente do período}
      {1.557 processos (879 da 5ª Turma e 678 da 6ª). Em 1.494 deles, quem recorreu foi a defesa.}
      {lida por um programa a partir do resumo oficial da decisão (a ``ementa''), sem leitura humana}

\begin{center}\small
\begin{tabular}{>{\ttfamily}lp{6.6cm}p{5cm}}
\toprule
\normalfont\bfseries Coluna & \bfseries O que diz & \bfseries Valores possíveis \\
\midrule
turma & qual turma julgou & 5ª ou 6ª \\
relator & o ministro responsável pelo processo & 11 nomes \\
classe & o tipo de recurso & HC, RHC, AREsp, REsp \\
objeto & o assunto principal do pedido & § 4º, busca ilegal, pena, \ldots \\
recorrente & quem recorreu & defesa ou Ministério Público \\
dispositivo & o que o tribunal decidiu & deu, negou, não conheceu, \ldots \\
resultado\_reu & o resultado \emph{para o réu} & favorável, parcial, desfavorável \\
barreiras & regras processuais usadas para não analisar o pedido & Súmula 7, Súmula 182, \ldots \\
marcadores & expressões sobre a pessoa ou o lugar & ``vulgo'', facção, ``local conhecido'', \ldots \\
votacao & se a decisão foi unânime & unânime ou por maioria \\
relator\_vencido & se o relator perdeu a votação & sim ou não \\
\bottomrule
\end{tabular}
\end{center}

O programa que lê as ementas foi conferido contra os processos lidos por inteiro: acertou
o resultado para o réu em 32 de 33 casos da 5ª Turma, 29 de 30 da 6ª e 16 de 19 da base do
§~4º.

\paragraph{Tabela 2 -- Amostras lidas por inteiro.}
\ficha{\texttt{dados/amostra\_casos.csv} (5ª Turma) e \texttt{dados/amostra\_t6\_casos.csv} (6ª Turma)}
      {um réu}
      {5ª Turma: 42 réus (34 dentro do recorte). 6ª Turma: 43 réus (31 dentro do recorte). Mais 14 casos escolhidos à mão para contraste.}
      {40 processos sorteados por turma, com o mesmo número por relator; o texto completo foi lido com apoio de um programa de IA, seguindo um roteiro escrito, e parte foi conferida à mão}

As colunas descrevem os \textbf{fatos} (droga, quantidade, como foi a abordagem policial,
antecedentes) e \textbf{o que cada instância decidiu} (juiz, Tribunal de Justiça e STJ).

\paragraph{Tabela 3 -- Base do § 4º.}
\ficha{\texttt{dados/amostra\_p4\_casos.csv}}
      {um réu que pediu a redução do § 4º}
      {22 réus (19 processos) até agora, todos do relator Sebastião Reis. O plano prevê 99 processos de 10 relatores.}
      {todas as vitórias da defesa e 8 derrotas sorteadas por relator; texto completo lido com o mesmo roteiro}

Esta tabela foi feita para ter os fatos que a lei exige. Principais colunas:
droga, quantidade em gramas, primário, reincidente, maus antecedentes, atos infracionais
(quando adolescente), processos em andamento, condenação por associação, indícios de
dedicação (balança, celular, investigação prévia\ldots), o que o juiz, o TJ e o STJ decidiram
sobre a redução e por quê.

\paragraph{Tabela 4 -- Entrada do juiz virtual.}
\ficha{\texttt{prolog/fatos\_casos.pl} (gerada pelo programa \texttt{scripts/gerar\_fatos.py})}
      {um réu em que a redução do § 4º foi discutida}
      {63 réus: 22 da base do § 4º, 17 da amostra da 5ª Turma, 13 da 6ª e 11 dos casos de contraste}
      {as Tabelas 2 e 3, convertidas para o formato que o juiz virtual lê}

A Tabela~\ref{tab:cobertura} mostra um ponto importante: \textbf{muitos dos fatos que a lei
exige não aparecem nas decisões}. A primariedade, por exemplo, só é informada em 28 dos 63
réus.

\begin{table}[h]
\centering\small
\caption{Em quantos dos 63 réus cada fato aparece no texto das decisões.}
\label{tab:cobertura}
\begin{tabular}{lr}
\toprule
Fato & Réus em que aparece \\
\midrule
Tipo de droga & 51 \\
Algum indício de dedicação (balança, celular, arma\ldots) & 49 \\
Quantidade em gramas & 40 \\
Se é primário & 28 \\
Antecedentes & 19 \\
Data do fato (completa) & 14 \\
Se é reincidente & 10 \\
Atos infracionais & 8 \\
Processos em andamento & 4 \\
\bottomrule
\end{tabular}
\end{table}

\subsubsection*{O juiz virtual (árvore A)}

O juiz virtual é um programa em Prolog (\texttt{prolog/arvore\_a.pl}), uma linguagem em que se
escrevem regras do tipo ``se \ldots\ então \ldots''. Para cada réu, ele faz estas perguntas:
\begin{enumerate}
  \item O fato aconteceu depois da Lei de Drogas de 2006? Se não, a redução não existe.
  \item O réu é primário?
  \item Tem bons antecedentes?
  \item Dedica-se a atividades criminosas?
  \item Faz parte de organização criminosa?
  \item Se passou em todas: qual fração aplicar (de 1/6 a 2/3)?
  \item Qual é a pena final, o regime e se dá para trocar a prisão por pena alternativa?
\end{enumerate}
As perguntas 2, 3 e 5 dependem só dos fatos. As perguntas 4 e 6 dependem de
\textbf{escolhas}, porque a lei não diz como respondê-las. No programa, cada escolha virou
um ajuste explícito, e um conjunto de ajustes forma uma \textbf{leitura} da lei.
Criamos três leituras (Tabela~\ref{tab:perfis}). Elas foram inspiradas em decisões reais
que estudamos, mas não pretendem representar nenhum ministro.

\begin{table}[h]
\centering\small
\caption{As escolhas que a lei não faz, e as três leituras usadas.}
\label{tab:perfis}
\begin{adjustbox}{max width=\textwidth}
\begin{tabular}{lccc}
\toprule
Escolha & Garantista & Intermediária & Restritiva \\
\midrule
Muita droga indica dedicação ao crime? & não & sim & sim \\
Muita droga, sozinha, basta para negar? & não & não & sim \\
O que é ``muita'' maconha / cocaína & --- & 10 kg / 1 kg & 1 kg / 100 g \\
Atos infracionais de quando era adolescente contam? & não & até 3 anos antes & até 5 anos antes \\
Processos ainda em andamento contam? & não & não & sim \\
Condenação antiga (mais de 5 anos) é esquecida? & sim & não & não \\
Quais indícios de dedicação contam & arma, facção & + balança, celular, & + denúncias, ``local \\
 & & investigação, transporte & conhecido'', dinheiro \\
Quantos indícios bastam & 2 & 2 & 1 \\
\bottomrule
\end{tabular}
\end{adjustbox}
\end{table}

Há ainda uma escolha que vale para todas as leituras: \textbf{o que fazer quando um fato não
aparece no processo}. Pode-se decidir a favor do réu (``na dúvida, a favor do réu'') ou
contra ele (``cabe à defesa provar''). Testamos as duas opções.

Para isolar o efeito da redução, o juiz virtual parte sempre da pena mínima de 5 anos e só
aplica a redução e o aumento por circunstâncias especiais (como tráfico perto de escola).
Ele também explica cada decisão: quais condições falharam e quais indícios foram usados.

\subsubsection*{A árvore aprendida (árvore B)}

A árvore B não foi escrita por nós: um programa (\texttt{scikit-learn}) a \textbf{aprendeu}
olhando as decisões reais. Ela aprende uma sequência de perguntas de sim ou não que melhor
separa os casos em que a defesa ganhou dos casos em que perdeu. Depois, a árvore é
convertida em regras Prolog (\texttt{prolog/arvore\_b.pl}).
\begin{itemize}
  \item \textbf{Linhas usadas:} os 1.494 processos da Tabela 1 em que a defesa recorreu.
        Desses, 40 foram vitórias da defesa.
  \item \textbf{O que ela tenta adivinhar (coluna-alvo):} se a defesa venceu ou perdeu.
  \item \textbf{O que ela pode olhar (colunas de entrada):} turma, relator, tipo de recurso,
        assunto e os marcadores da ementa. Ou seja, só o que se sabia antes do julgamento.
  \item \textbf{O que ela não pode olhar:} as barreiras processuais, a votação e se o relator
        perdeu, porque essas informações já fazem parte da decisão. Usá-las seria ``colar''
        a resposta.
  \item \textbf{Tamanho:} no máximo 4 perguntas seguidas, e cada resposta final precisa valer
        para pelo menos 10 processos.
  \item \textbf{Duas versões:} uma pode usar o relator e a turma; a outra não.
\end{itemize}
O objetivo não é acertar muito. Como a defesa perde em 97\% dos casos, um programa que
sempre diz ``perde'' acerta 97\% sem aprender nada. O que interessa é \textbf{quais
perguntas} ela aprende a fazer.

\subsubsection*{Como verificamos os programas}

O juiz virtual passou por 26 testes com réus inventados, um para cada regra (por exemplo:
``réu reincidente nunca recebe a redução'' e ``um fato de 2005 não tem redução''). A árvore B
foi testada para garantir que as regras em Prolog decidem exatamente como o modelo
aprendido, nos 1.494 processos.

% =====================================================================
\subsection{Experimentos Computacionais}
% =====================================================================

\subsubsection*{E1 -- A mesma lei, leituras diferentes}

\textbf{Pergunta:} se só as escolhas mudam, o resultado muda?
\textbf{Como:} os mesmos 63 réus foram julgados pelas três leituras, com o fato faltante
contando a favor e contra o réu.

\begin{table}[h]
\centering\small
\caption{E1 -- Os mesmos 71 réus julgados pelo juiz virtual com leituras diferentes da lei.}
\label{tab:e1}
\begin{adjustbox}{max width=\textwidth}
\begin{tabular}{llrrrr}
\toprule
Leitura (perfil) & Dado que falta conta & Réus com redução & Pena média (meses) & Regime aberto & Regime fechado \\
\midrule
garantista & a favor do réu & 63/71 & 34.7 & 52 & 3 \\
intermediário & a favor do réu & 30/71 & 47.6 & 26 & 3 \\
restritivo & a favor do réu & 8/71 & 60.1 & 3 & 3 \\
garantista & contra o réu & 16/71 & 56.9 & 11 & 3 \\
intermediário & contra o réu & 6/71 & 60.0 & 4 & 3 \\
restritivo & contra o réu & 0/71 & 62.3 & 0 & 3 \\
\bottomrule
\end{tabular}
\end{adjustbox}
\end{table}


\textbf{Resultado:} com o fato faltante a favor do réu, a leitura garantista dá a redução a
57 réus; a restritiva, a 7. As três leituras discordam em 50 dos 63 réus. Para o mesmo réu,
a pena varia em média 26 meses entre a leitura mais branda e a mais dura (até 47 meses).
Quando o fato faltante conta contra o réu, as reduções quase desaparecem (13, 4 e 0),
porque a primariedade, por exemplo, não aparece em 35 dos 63 processos.

Para saber \textbf{qual escolha pesa mais}, partimos da leitura intermediária e trocamos uma
escolha de cada vez:

\begin{table}[h]
\centering\small
\caption{E1 -- Trocando uma escolha de cada vez (63 réus).}
\label{tab:e1b}
\begin{adjustbox}{max width=\textwidth}
\begin{tabular}{p{5.6cm}lp{1.9cm}p{1.9cm}p{2.4cm}}
\toprule
Escolha trocada (partindo da leitura intermediária) & Na direção & Réus que mudam de decisão & Réus que mudam de fração & Réus com redução (antes $\to$ depois) \\
\midrule
indícios aceitos: só arma e facção & garantista & 24 & 0 & 26 $\to$ 50 \\
basta 1 indicador & restritivo & 13 & 0 & 26 $\to$ 13 \\
indícios aceitos: lista ampla & restritivo & 10 & 0 & 26 $\to$ 16 \\
quantidade não indica dedicação & garantista & 9 & 0 & 26 $\to$ 35 \\
limites de quantidade menores & restritivo & 3 & 0 & 26 $\to$ 23 \\
quantidade sozinha basta & restritivo & 2 & 0 & 26 $\to$ 24 \\
esquecimento de antecedentes (5 anos) & garantista & 2 & 0 & 26 $\to$ 28 \\
atos infracionais até 5 anos & restritivo & 1 & 0 & 26 $\to$ 25 \\
ato infracional sem data conta & restritivo & 1 & 0 & 26 $\to$ 25 \\
ações em curso contam (contra o Tema 1.139) & restritivo & 1 & 0 & 26 $\to$ 25 \\
faixas de fração restritivas & restritivo & 0 & 10 & 26 $\to$ 26 \\
faixas de fração garantistas & garantista & 0 & 3 & 26 $\to$ 26 \\
atos infracionais não contam & garantista & 0 & 0 & 26 $\to$ 26 \\
\bottomrule
\end{tabular}
\end{adjustbox}
\end{table}


As escolhas que mais mudam o resultado são \textbf{quais indícios contam} (até 24 réus) e
\textbf{quantos indícios bastam} (13). A regra sobre processos em andamento (Tema 1.139 do
STJ) mexe em poucos casos aqui porque esse fato quase nunca aparece nos textos.

\subsubsection*{E2 -- O mesmo processo, dois julgamentos}

\textbf{Pergunta:} as leituras do juiz virtual correspondem a decisões que realmente
aconteceram?
\textbf{Como:} em quatro processos da base do § 4º, o relator deu a redução sozinho, o
Ministério Público recorreu e a Turma tirou a redução. Comparamos as duas decisões com
as três leituras.

\begin{table}[h]
\centering\small
\caption{E2 -- O relator deu a redução sozinho e a Turma tirou (relator: Sebastião Reis; recurso do MP).}
\label{tab:e2}
\begin{adjustbox}{max width=\textwidth}
\begin{tabular}{llllllll}
\toprule
Processo & Droga apreendida & Relator (sozinho) & Turma & Votação & Garant. & Interm. & Restrit. \\
\midrule
AREsp 2946226 & 432 kg (maconha) & concede 1/6 & negado & maioria & concede 1/6 & nega & nega \\
HC 1017077 & 2 kg (cocaína) & concede (ofício) 2/3 & reduz a 1/6 & maioria & concede 1/3 & nega & nega \\
HC 1017446 & 17 kg (cocaína) & concede 2/3 & negado & unânime & concede 1/6 & nega & nega \\
HC 1060758 & 308 kg (maconha) & concede (ofício) 1/6 & negado & maioria & concede 1/6 & nega & nega \\
\bottomrule
\end{tabular}
\end{adjustbox}
\end{table}


\textbf{Resultado:} nos quatro casos, a decisão do relator coincide com a leitura
garantista (que dá a redução), e a decisão da Turma coincide com as leituras intermediária
e restritiva (que negam). Os fatos e a lei eram os mesmos; mudou a leitura. Em três dos
quatro casos a Turma decidiu por maioria, com o relator vencido.

\subsubsection*{E3 -- O que a árvore aprendida aprende}

\textbf{Pergunta:} quando um programa aprende com as decisões reais, ele aprende a lei ou
aprende quem julga?
\textbf{Como:} as duas versões da árvore B foram testadas em processos que não viram no
treino.

\begin{table}[h]
\centering\small
\caption{E3 -- Árvore aprendida: 5.323 recursos da defesa, dos quais 308 vitórias. Testada em casos que ela não viu no treino.}
\label{tab:e3}
\begin{adjustbox}{max width=\textwidth}
\begin{tabular}{p{3.4cm}lp{1.6cm}p{1.6cm}p{1.9cm}p{1.4cm}p{1.3cm}}
\toprule
Versão & 1ª pergunta & Acerto equilibrado (\%) & Vitórias achadas (\%) & Quando diz vence, acerta (\%) & Vitórias achadas & Alarmes falsos \\
\midrule
com relator & relator = DANIELA TEIXEIRA & 75 & 65 & 21 & 200/308 & 765 \\
sem relator & é REsp & 67 & 71 & 11 & 218/308 & 1840 \\
sem relator, testada num relator que não viu &  & 59 & 52 & 9 & 159/308 & 1674 \\
\bottomrule
\end{tabular}
\end{adjustbox}
\end{table}


As regras aprendidas pela versão com relator são estas:
\begin{itemize}
  \item Se o recurso é um REsp: a defesa vence se o relator é Sebastião Reis, Maria Marluce
        ou Og Fernandes; senão, perde.
  \item Se não é REsp: a defesa vence se o relator é Og Fernandes e o assunto é busca
        ilegal; senão, perde.
\end{itemize}
Sem acesso ao relator, a árvore passa a usar expressões da ementa (``fuga'', ``denúncia
anônima'', ``local conhecido'') e piora, principalmente quando é testada em um relator
que não viu.

\textbf{Resultado:} depois do tipo de recurso, todas as perguntas da árvore são sobre
\textbf{quem é o relator}, e nenhuma sobre os fatos que a lei manda considerar (os fatos nem
aparecem na Tabela 1). Se trocarmos só o nome do relator, mantendo tudo o mais, a resposta
da árvore muda em 292 dos 1.494 processos. Atenção: a árvore acha 27 das 40 vitórias, mas
dá 171 alarmes falsos; quando diz ``vence'', acerta só 14\% das vezes. Ela não serve como
juiz, e esse é justamente o ponto: o que ela aprendeu foi o padrão de quem decide.

\subsubsection*{E4 -- O juiz virtual e as decisões reais}

\textbf{Pergunta:} qual leitura se parece mais com o que os tribunais decidiram?
\textbf{Como:} contamos em quantos réus a decisão do juiz virtual sobre a redução é igual
à do Tribunal de Justiça e à do STJ.

\begin{table}[h]
\centering\small
\caption{E4 -- Juiz virtual $\times$ decisão real sobre o § 4º. O TJ concedeu em 11 de 68 réus; o STJ, em 11 de 60.}
\label{tab:e4}
\begin{adjustbox}{max width=\textwidth}
\begin{tabular}{llrr}
\toprule
Leitura (perfil) & Dado que falta conta & Igual ao TJ & Igual ao STJ \\
\midrule
garantista & a favor do réu & 19/68 & 19/60 \\
intermediário & a favor do réu & 40/68 & 39/60 \\
restritivo & a favor do réu & 57/68 & 47/60 \\
garantista & contra o réu & 45/68 & 44/60 \\
intermediário & contra o réu & 51/68 & 50/60 \\
restritivo & contra o réu & 57/68 & 49/60 \\
\bottomrule
\end{tabular}
\end{adjustbox}
\end{table}


\textbf{Resultado:} a leitura restritiva é a que mais se parece com os tribunais (41 de 52
no STJ). Mas isso precisa ser lido com cuidado: o STJ negou a redução em 42 dos 52 réus, então
um programa que sempre negasse acertaria 42. A concordância alta mostra que os tribunais
negam quase sempre, não que a leitura restritiva seja a ``correta''. A leitura garantista,
que segue a lei de forma mais literal, é a que mais discorda.

\subsubsection*{Limitações}

\begin{itemize}
  \item Só há decisões colegiadas. Quando a defesa ganha com o relator sozinho e o Ministério
        Público não recorre, o caso não aparece nos arquivos do STJ. Por isso, as taxas de
        vitória podem estar subestimadas.
  \item A leitura dos textos completos foi feita com apoio de IA, com conferência manual
        de parte dos casos. \todo{número de casos conferidos à mão}
  \item A base do § 4º ainda tem só um relator. \todo{completar com os outros 9 relatores}
  \item As três leituras são ilustrativas. O ponto não é qual delas é a certa, e sim que o
        programa precisa de alguma, e cada uma dá resultados diferentes.
\end{itemize}

\end{document}
